%%%%%%%%%%%%%%%%%%%%%%%%%%%%%%%%%%%%%%%%%%%%%%%%%%%%%%%%%%%%%%%%%%%%%%%%%%
%%%%%% Ajit Kumar Sahoo, School of Computer and Information Sciences,
%%%%%% University of Hyderabad, Hyderabad, India-500046
%%%%%% Email: sahooajitkumar85@gmail.com
%%%%%% https://ajitsahoocs.github.io/
%%%%%%%%%%%%%%%%%%%%%%%%%%%%%%%%%%%%%%%%%%%%%%%%%%%%%%%%%%%%%%%%%%%%%%%%%%%
\documentclass{beamer}
\hypersetup{pdfpagemode=FullScreen} %full screen mode
\setbeamertemplate{navigation symbols}{}
\usepackage[english]{babel} %designed for typesetting lebels(addr lebel,sticky lebels,etc)
\usepackage[latin1]{inputenc} %Accept different input encodings
\usepackage{amsmath,amssymb} 
\usepackage[T1]{fontenc} %for font encoding
\usepackage{times} % use times font instead of default
\usepackage{curves} %for drawing curves
\usepackage{verbatim} %paragraph making environment 
\usepackage{multimedia} %for multimedia like animation,movie etc...
\usepackage{mathptmx} % font style
\usepackage{graphicx} % Allows including images
\usepackage{booktabs} % Allows the use of \toprule, \midrule and \bottomrule in tables
\usepackage{hyperref}
\usepackage{xcolor}

\usepackage{algorithm,algorithmic}
\renewcommand{\algorithmicrequire}{\textbf{Input:}}
\renewcommand{\algorithmicensure}{\textbf{Output:}}
\usepackage{lipsum}
\setbeamertemplate{caption}[numbered] % For numbering figures



\mode<presentation> {

% The Beamer class comes with a number of default slide themes
% which change the colors and layouts of slides. Below this is a list
% of all the themes, uncomment each in turn to see what they look like.

%\usetheme{default}
%\usetheme{AnnArbor}
%\usetheme{Antibes}
%\usetheme{Bergen}
%\usetheme{Berkeley}
%\usetheme{Berlin}
%\usetheme{Boadilla}
%\usetheme{CambridgeUS}
%\usetheme{Copenhagen}
%\usetheme{Darmstadt}
%\usetheme{Dresden}
%\usetheme{Frankfurt}
%\usetheme{Goettingen}
%\usetheme{Hannover}
%\usetheme{Ilmenau}
%\usetheme{JuanLesPins}
%\usetheme{Luebeck}
\usetheme{Madrid}
%\usetheme{Malmoe}
%\usetheme{Marburg}
%\usetheme{Montpellier}
%\usetheme{PaloAlto}
%\usetheme{Pittsburgh}
%\usetheme{Rochester}
%\usetheme{Singapore}
%\usetheme{Szeged}
%\usetheme{Warsaw}

% As well as themes, the Beamer class has a number of color themes
% for any slide theme. Uncomment each of these in turn to see how it
% changes the colors of your current slide theme.

%\usecolortheme{albatross}
%\usecolortheme{beaver}%this one also
%\usecolortheme{beetle}
%\usecolortheme{crane} % also use this one
%\usecolortheme{dolphin}
%\usecolortheme{dove}
%\usecolortheme{fly}
%\usecolortheme{lily}
%\usecolortheme{orchid}
%\usecolortheme{rose}
%\usecolortheme{seagull}
%\usecolortheme{seahorse}
\usecolortheme{whale} % Best one
%\usecolortheme{wolverine} %use can use this also

%\setbeamertemplate{footline} % To remove the footer line in all slides uncomment this line
%\setbeamertemplate{footline}[page number] % To replace the footer line in all slides with a simple slide count uncomment this line

%\setbeamertemplate{navigation symbols}{} % To remove the navigation symbols from the bottom of all slides uncomment this line
}

\usepackage{graphicx} % Allows including images
\usepackage{booktabs} % Allows the use of \toprule, \midrule and \bottomrule in tables
\setbeamercovered{transparent}
\setbeamertemplate{bibliography item}[text]
\setbeamertemplate{theorems}[numbered]
\setbeamerfont{title}{size=\Large}%\miniscule,command,tiny, scriptsize,footnotesize,small,normalsize,large,Large,LARGE,huge,Huge,HUGE
\setbeamerfont{date}{size=\tiny}%{\fontsize{40}{48} \selectfont Text}

\setbeamertemplate{itemize items}[ball] % if you want a ball
\setbeamertemplate{itemize subitem}[circle] % if you wnat a circle
\setbeamertemplate{itemize subsubitem}[triangle] % if you want a triangle


%------------------customized frame----------------------------
\newcounter{cont}

\makeatletter
%allowframebreaks numbering in the title
\setbeamertemplate{frametitle continuation}{%
   % \setcounter{cont}{\beamer@endpageofframe}%
    %\addtocounter{cont}{1}%
   % \addtocounter{cont}{-\beamer@startpageofframe}%
   % (\insertcontinuationcount/\arabic{cont})%
}
\makeatother
%-----------------------customized-----------------------------




%----------------------------------------------------------------------------------------
%	TITLE PAGE
%----------------------------------------------------------------------------------------

\title[WiFi sensing]{Title Title Title Title Title Title Title Title Title Title Title Title Title } % The short title appears at the bottom of every slide, the full title is only on the title page

%\author[AKS]{Ajit Kumar Sahoo \texorpdfstring{\scriptsize Regd No: 16MCPC03}{}}
\author[AKS]{\texorpdfstring{Name\newline{\small Regd No:}}{Author}}

\institute[UoH] % Your institution as it will appear on the bottom of every slide, may be shorthand to save space
{

{\small Supervisor}\\
\medskip
{\normalsize Name}\\
%\begin{center}
%\includegraphics[width=0.1\textwidth]{uoh.png}
%\end{center}
\vspace{2cm}
\medskip
School of Computer and Information Sciences\\
University of Hyderabad% Your institution for the title page

%\medskip
%\textit{john@smith.com} % Your email address
}
%This will place the image at position "30 right/left and 120 up/down" relative to the top left corner of the current page.
\titlegraphic{%
  \begin{picture}(0,0)
    \put(30,125){\makebox(0,0)[rt]{\includegraphics[width=2cm]{uoh.png}}}
  \end{picture}}
 \date{\today} % Date, can be changed to a custom date
\begin{document}

\begin{frame}
\titlepage % Print the title page as the first slide
\end{frame}
\begin{frame}
  \frametitle{Contents}
 % \tableofcontents[pausesections,shaded]
 \tableofcontents
\end{frame}
% TABLE OF CONTENTS AT BEGIN OF EACH SECTION
\AtBeginSection[]{
  \begin{frame}<beamer>
    \frametitle{Current Section}
   \tableofcontents[currentsection]
  \end{frame}}
%----------------------------------------------------------------------------------------
%	PRESENTATION SLIDES
%----------------------------------------------------------------------------------------


%-------------------------------------------------------------------------------
%-------------------------------INTRODUCTION--------------------------------
%-----------------------------------------------------------


\section{Introduction} 
\begin{frame}[allowframebreaks]
\frametitle{Introduction}

   \begin{itemize}
       
       \item statement 1
       \item statement 2
       
       \item statement 3
   \end{itemize}
  \begin{figure} %htbp! %https://www1.cmc.edu/pages/faculty/aaksoy/latex/latexfive.html
	\includegraphics[trim=0cm 0cm 0cm 0cm,clip,width=0.5\textwidth]{images/figure1.jpg}
	\caption{Block diagram IoT Edge layer }
	\label{figure1}
\end{figure}
\end{frame}


\subsection{Motivation}
\begin{frame}[allowframebreaks]
\frametitle{Motivation}
\end{frame}

\subsection{Research Gaps}
\begin{frame}[allowframebreaks]
\frametitle{Research Gaps}
\end{frame}


\subsection{Objective}
\begin{frame}[allowframebreaks]
\frametitle{Objective}
\end{frame}




%-------------------------------------------------------------------------------------------------------------------------------
\section{Background and Literature Survey}
\begin{frame}[allowframebreaks]
\frametitle{Background and Literature Survey}
\begin{itemize}
\item statement 1
\item statement 2
\item statement 3
\item statement 4
\end{itemize}


\begin{equation}
     y = Hx + n 
\end{equation}

\end{frame}






%------------------------------------------------------------------------------------------------------------------------------
\section{Methodology}
\begin{frame}[allowframebreaks]
\frametitle{Methodology}
\begin{itemize}
\item statement 1
\item statement 2
\item statement 3
\item statement 4\footnotemark. 
\end{itemize}
\footnotetext[1]{\textit{A. K. Sahoo, V. Kompally and S. K. Udgata, "Wi-Fi Sensing based Real-Time Activity Detection in Smart Home Environment," 2023 IEEE Applied Sensing Conference (APSCON), Bengaluru, India, 2023, pp. 1-3.}}
\end{frame}



\begin{frame}{Methodology}

\begin{algorithm}[H]
\begin{algorithmic}[1]
\REQUIRE input data
\ENSURE output data 
\FOR{$j=1$ to $N$}
\STATE Statement
\ENDFOR
\end{algorithmic}
\caption{pseudocode for algorithms}
\label{alg1}
\end{algorithm}
\end{frame}


\begin{frame}{Methodology}
\begin{definition} 
        This is a definition.
    \end{definition}
    
    \begin{theorem} 
        This is a theorem. 
    \end{theorem}
    
    \begin{lemma} 
        This is a proof idea.
    \end{lemma}


\end{frame}




\begin{frame}{Methodology}
    \begin{block}{Standard Block}
        This is a standard block.
    \end{block}
    \begin{alertblock}{Alert Message}
        This block presents alert message.
    \end{alertblock}
    \begin{exampleblock}{An example of typesetting tool}
        Example: MS Word, \LaTeX{}
    \end{exampleblock}
  
\end{frame}



\begin{frame}{Methodology}
    \begin{proof} 
        This is a proof. 
    \end{proof}
    
    \begin{corollary}
        This is a corollary
    \end{corollary}
    
    \begin{example}
        This is an example 
    \end{example}
    
\end{frame}
%-------------------------------------------------------------------------------------------------------------------------------
\section{Experiments}
\begin{frame}[allowframebreaks]
\frametitle{Experiments}
\begin{itemize}
\item statement 1
\item statement 2
\item statement 3
\item statement 4
\end{itemize}
\end{frame}
%--------------------------------------------------------------------------------------------------------------------------------
\section{Results}
\begin{frame}[allowframebreaks]
\frametitle{Results}
\begin{itemize}
\item statement 1
\item statement 2
\item statement 3
\item statement 4
\end{itemize}
\end{frame}

%-----------------------------------------------------------------------------------------
\section{Conclusion and future work}
\begin{frame}[allowframebreaks]
\frametitle{Conclusions and future work}
\begin{itemize}
\item statement 1
\item statement 2
\item statement 3
\item statement 4
\end{itemize}
\end{frame}

%--------------------------------------------------------------------------------------------------------------------------------
%\section{Information processing  issues}
%-----------------------------------------------
%\begin{frame}
%\frametitle{Information processing issues}
%\begin{itemize} 
%\item In collaborative processing
%\begin{itemize}
%\item Target detection, localization, communication, storage, querying...
%\end{itemize}
%\color{lightgray}
%\item  In networking
%\begin{itemize}
%\color{lightgray}
%\item  aggregation, routing, Data naming,...
%\end{itemize}
%\end{itemize}
%\end{frame}
%-----------------------------------------------
%\begin{frame}[fragile] % Need to use the fragile option when verbatim is used in the slide
%\frametitle{Citation}
%An example of the \verb|\cite| command to cite within the presentation:\\~
%Citation \cite{Souza} requires .
%\end{frame}
%------------------------------------------------
\section{References}
\begin{frame}[allowframebreaks]
\frametitle{References}
 \nocite{*}
 \bibliographystyle{unsrt}
\bibliography{References}
\end{frame}

%-----------------------------------------------------------------------------------------
\section{Acknowledgement}
\begin{frame} [allowframebreaks]
\frametitle{Acknowledgement}
\lipsum[1]


\end{frame}
%-----------------------------------------------------------------------------------------

\begin{frame}
\Huge{\centerline{Thank You}}
\end{frame}

%----------------------------------------------------------------------------------------

\end{document}
